% Baccalauréat malgache, série D, 2008 : énoncé.
\TCDBacAttribution{D}{2008}
\providecommand{\Courbe}{(\mathscr{C})}
\providecommand{\D}{(\mathscr{D})}
\providecommand{\Plan}{(\mathscr{P})}
\providecommand{\Ouv}{(O ; \vec{u} , \vec{v})}
%\providecommand{\Oij}{(O ; \vec{\imath} , \vec{\jmath})}
\providecommand{\cm}{~\mathrm{cm}}
\providecommand{\dx}{~\mathrm{d}x}
\providecommand{\Arg}[1]{\mathrm{Arg}\left( #1 \right)}

\noindent \underline{\textsc{Exercice I}} (5 points) 
      \begin{enumerate}
            \item R\'esoudre dans $\mathbb{C}$ l'\'equation : $z^2 - (1-3i) z - 4 = 0$.
            \item \begin{enumerate}
                        \item Dans le plan complexe $\Plan$, muni d'un rep\`ere orthonorm\'e $\Ouv$ d'unit\'e $2\cm$,
                              placer les points $A$, $B$ et $D$ d'affixes respectives $1$, $(-1-i)$ et $(2-2i)$ 
                              ainsi que le point $C$ milieu du segment $[BD]$ dont on pr\'ecisera l'affixe.
                        \item D\'eterminer l'ensemble $(\Delta)$ des points $M$ d'affixe $z$ tels que :
                              $\left| \frac{z +1 +i}{z -2 + 2i} \right| = 1$.
                        \item Montrer que $(\Delta)$ passe par $A$. Achever alors la construction de $(\Delta)$.
                  \end{enumerate}
            \item On consid\`ere la similitude plane directe $S$ d\'efinie par : 
                  $
                        \begin{multiform}
                              S(B) = B \\ S(A) = D \;.
                        \end{multiform}
                  $
                  Pr\'eciser les \'el\'ements caract\'eristiques de $S$.
      \end{enumerate}
%%%%%%%%%%%%%%%%%%%%%%%%%%%%%%%%%%%%%%%%%%%%%%%%%%%%%%%%%%%%%%%%%%%%%%%%%%%%%%%%%%%%%%%%%%%%%%%
%%%%%%%%%%%%%%%%%%%                   %%%%%%%%%%%%%%%%%%%                   %%%%%%%%%%%%%%%%%%%
%%%%%%%%%%%%%%%%%%%%%%%%%%%%%%%%%%%%%%%%%%%%%%%%%%%%%%%%%%%%%%%%%%%%%%%%%%%%%%%%%%%%%%%%%%%%%%%
\noindent \underline{\textsc{Exercice II}} (5 points) 
      \begin{description}
            \item[A. ]
                  On dispose d'un d\'e cubique bien \'equilibr\'e dont les faces sont num\'erot\'ees de 1 \`a 6.
                  \begin{enumerate}
                        \item On lance une fois le d\'e.
                              Calculer la probabilit\'e d'avoir un nombre strictement sup\'erieur \`a 2.
                        \item Maintenant, on lance deux fois de suite ce d\'e. 
                              Une \'eventualit\'e est un couple d'entiers naturels $(a,b)$.
                              On appelle $X$ la variable al\'eatoire qui, \`a chaque couple $(a,b)$ obtenu, 
                              associe le r\'eel $|a-b|$.
                              \begin{enumerate}
                                    \item D\'eterminer l'univers image de $X$.
                                    \item Etablir la loi de probabilit\'e de $X$.
                              \end{enumerate}
                  \end{enumerate}
            \item[B. ]
                  Lors d'un test, les notes obtenues par 4 candidats, aux \'epreuves de chant et de musique,
                  sont indiqu\'ees dans le tableau suivant :
                  \par\smallskip\noindent
                  \adjustbox{max width=\linewidth,center}{%
                  \begin{tabular}{|l|c|c|c|c|}
                        \hline
                        Musique ($x_i$) & $\alpha$ & 3 & 6 & 9       \\
                        \hline
                        Chant ($y_i$)   &     2    & 4 & 5 & $\beta$ \\
                        \hline
                  \end{tabular}%
                  }
                  \par\smallskip
                  \begin{enumerate}
                        \item On sait que le point moyen associ\'e \`a cette s\'erie statistique a pour coordonn\'ees 
                              $\bar{x} = 5$ et $\bar{y} = 4,5$ ;
                              d\'eterminer les notes $\alpha$ et $\beta$ respectivement obtenues par deux candidats diff\'erents
                              en musique et en chant.
                        \item D\'eterminer le coefficient de corr\'elation lin\'eaire de cette s\'erie.
                              Interpr\'eter le r\'esultat obtenu.
                        \item D\'eterminer l'\'equation de la droite de r\'egression de $y$ en $x$.
                  \end{enumerate}
      \end{description}
%%%%%%%%%%%%%%%%%%%%%%%%%%%%%%%%%%%%%%%%%%%%%%%%%%%%%%%%%%%%%%%%%%%%%%%%%%%%%%%%%%%%%%%%%%%%%%%
%%%%%%%%%%%%%%%%%%%                   %%%%%%%%%%%%%%%%%%%                   %%%%%%%%%%%%%%%%%%%
%%%%%%%%%%%%%%%%%%%%%%%%%%%%%%%%%%%%%%%%%%%%%%%%%%%%%%%%%%%%%%%%%%%%%%%%%%%%%%%%%%%%%%%%%%%%%%%
\noindent \underline{\textsc{Probleme}} (10 points) 
      \begin{description}
            \item[I. ] $g$ est la fonction num\'erique d\'efinie et continue sur $]0;+\infty[$ par $g(x) = x^2 + \ln x$.
                  \begin{enumerate}
                        \item Dresser le tableau de variation de $g$ sur $]0;+\infty[$.
                        \item Montrer qu'il existe un r\'eel unique $\alpha$ dans $]0;+\infty[$ tel que $g(\alpha) = 0$.
                        \item Etudier le signe de $g(x)$ suivant les valeurs de $x$.
                  \end{enumerate}
            \item[II. ]
                  On consid\`ere maintenant la fonction num\'erique $f$ d\'efinie sur $]0;+\infty[$ par 
                  $
                        f(x) = (x+1) - \left( \frac{1+\ln x}{x} \right) .
                  $
                  On d\'esigne par $\Courbe$ sa courbe repr\'esentative dans le plan $\Plan$ muni d'un rep\`ere orthonorm\'e
                  direct \Oij{} d'unit\'e $4\cm$.
                  \begin{enumerate}
                        \item \begin{enumerate}
                                    \item Montrer que $f'(x) = \frac{g(x)}{x^2}$ pour tout r\'eel $x$ de $]0;+\infty[$
                                          o\`u $f$ est la fonction d\'eriv\'ee de $f$.
                                    \item En d\'eduire que $f'(x)$ a le m\^eme signe que $g(x)$ 
                                          pour tout r\'eel $x$ de $]0;+\infty[$.
                              \end{enumerate}
                        \item \begin{enumerate}
                                    \item Calculer les limites de $f$ aux bornes de $]0;+\infty[$.
                                    \item Dresser le tableau de variation de $f$.
                              \end{enumerate}
                        \item \begin{enumerate}
                                    \item Montrer que la droite $(\Delta)$ d'\'equation $y = x+1$ est une asymptote oblique
                                          de $\Courbe$.
                                    \item D\'eterminer l'abscisse de $A$, point d'intersection de $\Courbe$ avec $(\Delta)$.
                              \end{enumerate}
                        \item Tracer $(\Delta)$ et $\Courbe$ dans le m\^eme rep\`ere \Oij .
                              (Pour construire, on prendra $\frac{1}{e} = 0,4$ ; $\alpha = 0,6$ et $f(\alpha) = 0,8$.)
                        \item Calculer l'int\'egrale $I$ d\'efinie par 
                              $\displaystyle{
                                    I = \int_{1}^{e} \frac{1+\ln x}{x} \dx
                              }$.
                  \end{enumerate}
      \end{description}








